\chapter{Addition with Carrying}\label{ch:g2:addition}

When the units of a sum pass nine, ten of them bundle into a new
ten --- and that little bundle, moved to the next column, is the
famous \emph{carry}. This chapter explains the carry with sticks,
then trains the written method.

\section{Why we carry}

\begin{example}[Bundling ten units]\label{ex:g2:addition:bundle}
$27 + 15$: put the units together, $7 + 5 = 12$ units. But ten
units make a ten! Bundle ten of them: that leaves $2$ loose units
and one new ten to carry. Now the tens: $2 + 1 + 1 = 4$ tens. So
$27 + 15 = 42$.
\end{example}

\begin{omfigure}
\begin{tikzpicture}[scale=0.75]
  % 12 units: ten of them tied into a bundle, two left loose
  \foreach \i in {0,...,9} \draw[omThm, very thick]
    (0.24*\i,0) -- (0.24*\i,1.7);
  \draw[omExo, thick, rounded corners=3pt]
    (-0.24,0.62) rectangle (2.4,1.05);
  \node[below, font=\small] at (1.08,-0.25) {$10$ bundled};
  \foreach \i in {0,1} \draw[omDef, very thick]
    (4.0+0.4*\i,0) -- (4.0+0.4*\i,1.7);
  \node[below, font=\small] at (4.2,-0.25) {$2$ left};
  \node[right, font=\small, align=left] at (5.5,0.85)
    {$7 + 5 = 12$: one ten and $2$ units.\\
     The bundle is the carry.};
\end{tikzpicture}

{\small The carry is a bundle: ten loose units become one ten, moved
to the tens column.}
\end{omfigure}

\section{The written method}

\begin{method}[Column addition]\label{met:g2:addition:column}
\begin{enumerate}
  \item Write the numbers one under the other, units under units,
        tens under tens;
  \item add the units; if the answer passes $9$, write the units
        digit and carry a small $1$ above the tens;
  \item add the tens (with the carry), then the hundreds;
  \item check with an estimate: round each number and add roughly.
\end{enumerate}
\end{method}

\begin{example}\label{ex:g2:addition:worked}
Compute $268 + 156$:
\[
\begin{array}{r}
{\scriptstyle 1}\ {\scriptstyle 1}\ \phantom{0} \\[-3pt]
2\,6\,8 \\
+\ 1\,5\,6 \\
\hline
4\,2\,4
\end{array}
\]
Units: $8 + 6 = 14$: write $4$, carry $1$. Tens:
$6 + 5 + 1 = 12$: write $2$, carry $1$. Hundreds:
$2 + 1 + 1 = 4$. Estimate check: about $270 + 160 = 430$, close to
$424$. \checkmark
\end{example}

\section{Adding in your head}

\begin{example}[Make a ten]\label{ex:g2:addition:mental}
For small sums, jump to a round number first:
\[
48 + 6 = 48 + 2 + 4 = 50 + 4 = 54 .
\]
Add tens and units separately:
$35 + 23 = (30 + 20) + (5 + 3) = 50 + 8 = 58$. And a near-hundred
trick: $99 + 46 = 100 + 46 - 1 = 145$.
\end{example}

\begin{example}[A little problem]\label{ex:g2:addition:problem}
``The library has $185$ books; $57$ new ones arrive. How many books
now?'' Adding: $185 + 57 = 242$ (units $5 + 7 = 12$, carry\,\dots).
The library has $242$ books.
\end{example}

\section{Exercises}

\begin{exercise}[$\star$]\label{exo:g2:addition:1}
Compute (no carry needed): $34 + 25$; \quad $126 + 43$; \quad
$507 + 82$.
\end{exercise}

\begin{exercise}[$\star$]\label{exo:g2:addition:2}
Compute with a carry: $47 + 8$; \quad $56 + 27$; \quad $65 + 19$.
\end{exercise}

\begin{exercise}[$\star$]\label{exo:g2:addition:3}
Compute in columns: $268 + 145$; \quad $374 + 158$; \quad
$429 + 293$.
\end{exercise}

\begin{exercise}[$\star$]\label{exo:g2:addition:4}
Compute in your head by making a ten: $38 + 5$; \quad $67 + 6$;
\quad $54 + 8$.
\end{exercise}

\begin{exercise}[$\star$]\label{exo:g2:addition:5}
Compute by adding tens and units separately: $42 + 36$; \quad
$53 + 24$; \quad $61 + 28$.
\end{exercise}

\begin{exercise}[$\star$]\label{exo:g2:addition:6}
Estimate first (round to the nearest ten or hundred), then compute
exactly: $58 + 34$; \quad $296 + 187$.
\end{exercise}

\begin{exercise}[$\star$]\label{exo:g2:addition:7}
``Ana has $128$ marbles, Tom has $95$. How many together?'' Write
the addition and the answer sentence.
\end{exercise}

\begin{exercise}[$\star$]\label{exo:g2:addition:8}
Complete the units hole: $3\,?\, + 27 = 64$ (find the missing units
digit of the first number).
\end{exercise}

\begin{exercise}[$\star$]\label{exo:g2:addition:9}
Two classes go to the show: $27$ and $28$ children, plus $6$
grown-ups. How many people in all? (Add the two classes first.)
\end{exercise}

\begin{exercise}[$\star\star$]\label{exo:g2:addition:10}
Find three numbers on the list $27$, $33$, $40$, $65$ that add up
to $100$. (Look for two whose units are ten friends first.)
\end{exercise}
